%% -*- eval: (flyspell-mode 1); -*-

\documentclass[11pt,onecolumn]{report}
\usepackage[utf8]{inputenc}
\usepackage[T1]{fontenc}

\usepackage[french]{babel}

\usepackage{graphicx}
\graphicspath{{./img/}}

\usepackage{geometry}
\geometry{hmargin=2.5cm,vmargin=2.8cm}

\usepackage{tikz}

\usetikzlibrary{arrows,positioning,fit}


\usepackage{amsmath}

\usepackage{floatrow}
\usepackage{minted}

\makeatletter
\newcommand{\minted@write@detok}[1]{%
  \immediate\write\FV@OutFile{\detokenize{#1}}}%

\newcommand{\minted@FVB@VerbatimOut}[1]{%
  \@bsphack
  \begingroup
    \FV@UseKeyValues
    \FV@DefineWhiteSpace
    \def\FV@Space{\space}%
    \FV@DefineTabOut
    %\def\FV@ProcessLine{\immediate\write\FV@OutFile}% %Old, non-Unicode version
    \let\FV@ProcessLine\minted@write@detok %Patch for Unicode
    \immediate\openout\FV@OutFile #1\relax
    \let\FV@FontScanPrep\relax
%% DG/SR modification begin - May. 18, 1998 (to avoid problems with ligatures)
    \let\@noligs\relax
%% DG/SR modification end
    \FV@Scan}
    \let\FVB@VerbatimOut\minted@FVB@VerbatimOut

\renewcommand\minted@savecode[1]{
  \immediate\openout\minted@code\jobname.pyg
  \immediate\write\minted@code{\expandafter\detokenize\expandafter{#1}}%
  \immediate\closeout\minted@code}
\makeatother


\usepackage[backend=biber,style=numeric,doi=false,isbn=false,hyperref=true,backref=true]{biblatex}
\addbibresource{art.bib}
\addbibresource{online.bib}
\addbibresource{ths.bib}

\DefineBibliographyStrings{french}{%
    backrefpage  = {cf p.}, % for single page number
    backrefpages = {cf pp.} % for multiple page numbers
}

\usepackage{hyperref}
\hypersetup{colorlinks=true}


\begin{document}

\begin{titlepage}

  \begin{center}
    \begin{figure}[!htbp]
      \begin{center}
        \begin{minipage}{0.4\linewidth}
          \center{\includegraphics[scale=0.7]{enseirb_inp.png}}
        \end{minipage}
        \hfill
        \begin{minipage}{0.4\linewidth}
          \center{\includegraphics[scale=0.05]{UniversiteBordeaux_CMJN-01.jpg}}
        \end{minipage}
      \end{center}
    \end{figure}
  \end{center}

  \vspace{0.5cm}
  \hrule
  \begin{flushleft}
    \Huge{\textit{Extraction automatisée de parallélisme de tâches, et expression dans le formalisme SYCL}}\\
  \end{flushleft}
  \begin{flushright}
    \huge\textbf{Rapport de stage}\\
  \end{flushright}
  \hrule

  \vspace{1cm}
  \center{\large {Alexandre Honorat}}\\
  \vspace{1cm}
  \large{\textbf{Encadrants :} Olivier Aumage, Denis Barthou}\\
  \large{\textbf{Tuteur :} Denis Barthou}
  \vspace{1cm}

  \center{\small{Filière Informatique, spécialité Parallélisme, Régulation et Calcul Distribué\\Master Réseaux et Systèmes Mobiles, parcours Calcul Haute Performance\\\today}}

  \vspace{1cm}

  \begin{center}
    \begin{figure}[!htbp]
      \begin{center}
        \begin{minipage}{0.4\linewidth}
          \center{\includegraphics[scale=0.3]{labri.jpeg}}
        \end{minipage}
        \hfill
        \begin{minipage}{0.4\linewidth}
          \center{\includegraphics[scale=0.3]{INRIA_CORPO_CMJN.png}}
        \end{minipage}
      \end{center}
    \end{figure}
  \end{center}

\end{titlepage}


\tableofcontents

\chapter*{Préambule}

\section*{Cadre}

Ce rapport concerne le stage de fin d'études que j'ai effectué chez INRIA Bordeaux (Institut National de Recherche en Informatique et Automatisme) au sein de l'équipe Storm. Le stage s'est déroulé du $1^{er}$ février $2015$ au $31$ juillet $2015$, encadré par Olivier Aumage et Denis Barthou. Olivier Aumage est chargé de recherche dans l'équipe Storm. Denis Barthou est le directeur de cette équipe, il est également professeur à l'Institut Polytechnique de Bordeaux (IPB). L'intitulé du stage était \og Extraction automatisée de parallélisme de tâches, et expression dans le formalisme SYCL \fg.

Du point de vue académique ce stage sera évalué à la fois par l'Enseirb-Matmeca dans le cadre du diplôme d'ingénieur en informatique avec spécialité PRCD (Parallélisme, Régulation et Calcul Distribué) -- et par l'Université de Bordeaux dans le cadre du master de recherche informatique RSM (Réseaux et Systèmes Mobiles, parcours Calcul Haute Performance). À ce titre mes tuteurs respectifs sont Denis Barthou et Paul Dorbec.

\section*{Remerciements}

Durant toute la durée du stage, Olivier Aumage et Denis Barthou ont répondu à mes nombreuses questions, ce pour quoi je les remercie beaucoup. De plus ils m'ont laissé une grande autonomie sur le code produit, et m'ont permis de voir des aspects très divers de la problématique de base. Par ailleurs Ronan Keryell nous a fortement aidé en nous fournissant la brique de base du stage, un prototype d'implantation de SYCL. Enfin Samuel Pitoiset a eu la patience de m'aider dans l'installation de certains compilateurs, et d'améliorer l'un d'entre eux pour qu'il fonctionne avec notre outil.


%% -*- eval: (flyspell-mode 1); -*-

\chapter{Introduction}
%\addcontentsline{toc}{chapter}{Introduction}

Un des aspects de la recherche informatique consiste à améliorer des programmes afin de rendre leur temps d'exécution plus rapide, il s'agit du \emph{calcul haute performance} (ou HPC, pour High Performance Computing). 
%Les performances -- en matière de temps d'exécution -- sont primordiales par exemple pour les codes de simulation numérique dédiés à la météorologie qui calculent les prévisions du lendemain, où bien entendu la simulation doit donc s'être terminée en moins d'une nuit. 
Dans ce cadre le stage a porté sur la génération automatique de certains de ces programmes améliorés. 
Les sections suivantes présentent brièvement quelques rappels sur les problématiques HPC, ainsi que l'objectif du stage : l'adaptation automatique d'une certaine catégorie de problèmes au matériel prévu pour le HPC. Le chapitre $2$ précise plus particulièrement ces problématiques et l'objectif tandis que le chapitre $3$ se concentre sur le travail réalisé. Le chapitre $4$ en propose une évaluation et est suivi par la conclusion. 


\section{Parallélisation de programmes}

Diverses techniques existent pour améliorer les performances d'un programme et plus particulièrement de son temps d'exécution. Deux aspects peuvent entrer en compte : les algorithmes utilisés dans le programme ainsi que le matériel qui l'exécute. Le facteur limitant à améliorer en priorité est alors le matériel, car il influe sur la nature des algorithmes pouvant être utilisés.

La solution la plus simple pour améliorer les performances consisterait à se baser uniquement sur la loi de Moore, c'est-à-dire sur le fait que les capacités du matériel sur lequel le programme est exécuté sont doublées tous les ans. Or cette loi exponentielle n'est maintenant plus valable puisque la fréquence d'horloge des processeurs modernes stagne depuis quelques années à une valeur proche de $3$ GHz. 
%C'est cette fréquence d'horloge qui impose au matériel (si celui-ci est constitué de cet unique processeur) le nombre d'opérations qu'il va pouvoir effectuer en une seconde. Puisque la fréquence n'augmente plus, le nombre d'opérations par seconde pour ce type de matériel ne peut plus augmenter non plus, aucun gain de rapidité d'exécution n'est alors possible.

%La première solution s'appuyait ainsi sur l'augmentation de la fréquence, c'est-à-dire de la rapidité de chaque opération sachant qu'une seule est faite à la fois. Au lieu d'augmenter la rapidité des opérations, 
Une solution est d'augmenter le nombre d'opérations effectuées \emph{en même temps} par un processeur, ce qui est l'introduction du parallélisme dans le matériel.
%Celui-ci peut s'effectuer de plusieurs manières : à l'intérieur et à l'extérieur du processeur. À l'intérieur,
Il s'agit soit de subdiviser les instructions en  micro-instructions afin de créer un pipeline (ainsi une instruction peut commencer immédiatement après la fin de la première micro-instruction de l'instruction précédente), soit de dupliquer certains composants de l'architecture afin d'exécuter plusieurs instructions sur à la fois, soit enfin l'application de la même instruction sur plusieurs données à la fois. Les processeurs modernes intègrent tous ces trois technologies : le pipeline d'instructions, l'architecture nommée \emph{superscalaire}, et les calculs dits \emph{vectoriels}.
%À l'extérieur il s'agit de connecter plusieurs processeurs les uns avec les autres, tous participant aux calculs nécessités par le programme de base. Différentes manières de relier les processeurs existent, ils peuvent alors être appelés des cœurs ; dans toute la suite du rapport, un processeur est synonyme d'un ensemble de cœurs connectés entre eux par de la mémoire cache hiérarchique en accès uniforme\footnote{Cette notion sera explicitée au point \ref{sec:stencil_base}.}.
Les processeurs peuvent eux-mêmes être répliqués dans leur ensemble, soit sur une même surface de silicium (processeurs multicœurs), soit distinctement et alors reliés par un bus (multiprocesseurs), soit encore par une combinaison de ces deux techniques.
Enfin en dehors des processeurs classiques -- ou \emph{centraux}, abrégés en CPU -- existent aussi des processeurs graphiques -- dénommés GPU --, qui contiennent beaucoup de cœurs mais exécutent obligatoirement tous la même instruction en même temps. Bien sûr les deux types peuvent coexister sur une même machine.

La nature du matériel (et non seulement ses capacités d'ordre physique comme la fréquence d'horloge) nécessite une conception des algorithmes en adéquation (par exemple pour utiliser le côté superscalaire). 
%cela induit des modifications dans la façon de transcrire les algorithmes (par exemple pour utiliser le côté superscalaire), ainsi que dans les algorithmes eux-mêmes qui ne sont parfois plus du tout adaptés (notamment à cause des communications entre les processeurs, inexistantes auparavant). 
Écrire un programme destiné à une machine parallèle impose donc d'avoir des algorithmes souvent très spécifiques à la configuration de la machine d'exécution, ce qui est d'autant plus complexe avec les architectures des machines actuelles de plus en plus hétérogènes, et peut poser des problèmes de portabilité des applications.

\section{\'Ecriture de code parallèle}

%Les algorithmes étant souvent spécifiques au type de parallélisme et à l'architecture de la machine cible, les langages permettant d'écrire des programmes parallèles automatisent difficilement la génération de codes parfaitement adaptés à la machine cible. En pratique la plupart des langages modernes permettent l'expression du parallélisme sans le découvrir eux-même ; certains outils aident à écrire du code parallèle mais ils sont souvent expérimentaux et nécessitent alors que le code automatiquement produit soit vérifié avant d'être compilé\footnote{Lors de la transcription d'un problème en informatique, il est d'abord nécessaire de transcrire le problème vers un langage informatique (le code) -- compréhensible par les humains -- qui est alors \emph{compilé} par un programme préexistant (le compilateur) qui le transforme en fichier binaire exécutable -- compréhensible par \emph{la} machine.}.

Le procédé classique consiste à premièrement décrire le problème de manière algorithmique, deuxièmement à le transcrire dans un langage de programmation donnée -- et utiliser les outils natifs du langage ou alors les bibliothèques spécifiques pour le parallélisme -- et troisièmement à compiler le code écrit à l'étape précédente grâce à différents outils internes au compilateur. Les deux premières étapes sont réalisées par le développeur, à qui il convient de faciliter la tâche. 

Les difficultés inhérentes à la production d'un programme parallèle sont de plusieurs natures, notamment l'identification des sections parallèles et de celles obligatoirement séquentielles, ainsi que la gestion des communication et l'utilisation de toutes les ressources disponibles sur la matériel. L'identification des régions parallélisables ou non dépend de l'algorithme utilisé pour résoudre le problème voulu. Un programme bien parallélisé réduit au maximum le nombre d'opérations qui sont exécutées séquentiellement, afin que la majeure partie du temps d'exécution soit passée dans des sections parallèles où toute la puissance de la machine parallèle est utilisée. Par ailleurs il est essentiel de gérer les communications entre les différents processeurs afin que des calculs puissent être effectués pendant les communications ; on parle alors du recouvrement des communications. Enfin il convient d'utiliser le plus de ressources possibles sans que cette utilisation n'entraîne d'effort trop important, par exemple si de nombreuses parties de l'algorithme nécessitent une synchronisation de tous les processeurs -- ce qui est très coûteux en temps.

Actuellement c'est au programmeur qu'incombent presque toutes les tâches de parallélisation\footnote{En ce qui concerne les capacités des compilateurs, \textsf{gcc} comprend les directives de la bibliothèque \textsf{OpenMP}, mais ne les trouve pas automatiquement (bien qu'elles soient souvent très simples) et c'est donc au programmeur de les écrire. En revanche les compilateurs modernes sont capables d'exploiter eux-mêmes la caractère superscalaire d'un processeur, mais pas toujours de manière optimale. De même certains calculs vectoriels peuvent être détectés automatiquement, mais ils sont le plus souvent déclenchés via des fonctions spécifiques appelées par le développeur.}. %, or cela nécessite beaucoup de compétences, de temps, et de tests avant de parvenir à un résultat acceptable. S'il n'est pas toujours possible de découvrir le parallélisme automatiquement, il n'est pas non plus toujours facile de savoir où le décrire. Alors que 
Certaines options sont très proches du processeur (comme les capacités vectorielles) et doivent être gérées de manière assez fines, d'autres peuvent être prises en charge à un plus haut-niveau comme le lancement de plusieurs procédures distinctes en parallèle (a priori chacune des procédures étant alors associée à un processeur libre). %Cela amène à une hiérarchisation du parallélisme, ce qui complique encore plus la tâche du programmeur.
Par ailleurs certains paramètres du programme peuvent êtres calculés lors de la compilation -- paramètres statiques -- et d'autres uniquement lors de l'exécution -- paramètres dynamiques. Or pour améliorer la rapidité du programme, il est souhaitable de limiter les calculs lors de l'exécution. Lorsque le programmeur écrit un code en HPC, il espère donc avoir le plus possible de paramètres statiques, selon les possibilités du compilateur et du langage. 

Une des problématiques du parallélisme est donc de faciliter l'écriture de programme, soit en découvrant automatiquement le parallélisme -- ce qui est plutôt difficile -- , soit en fournissant au programmeur des outils adaptés -- ce qui est un peu plus facile.

\section{Choix du niveau d'expression du parallélisme}

%Pour écrire un programme informatique résolvant un problème donné, il est nécessaire de passer par de nombreuses étapes intermédiaires qui sont tout autant de stades où le parallélisme peut être décrit. De très nombreuses possibilités existent et nous ne les présenteront pas en intégralité ici. 

Compte-tenu de la complexité de l'écriture de code parallèle, des outils spécifiques ont été mis en place afin de faciliter le développement de tels programmes. Or ces outils -- langages, bibliothèques, ou compilateur -- sont rares et souvent trop spécifiques ou au contraire trop généraux. L'objectif du stage est donc de développer un nouvel outil adapté à un certain type de problèmes pour lequel peu d'outils existent et ne sont pas adaptés à la configuration hétérogène des machines actuelles.

Introduire un nouveau langage dédié au parallélisme est complexe car il faut non seulement développer le compilateur et les bibliothèques standards associés , mais aussi le rendre simple à utiliser et si possible pas trop éloigné des autres paradigmes dominants afin de faciliter son appréhension. D'un autre côté garder les langages existants implique de modifier les compilateurs pour qu'ils découvrent automatiquement les possibilités de parallélisation, ce qui est difficile à mettre en œuvre -- à cause de la complexité même de trouver le parallélisme, et à cause de la complexité des compilateurs modernes.
%Les deux solutions existent actuellement avec un succès mitigé\footnote{Ce constat est difficile à vérifier, mais à titre d'exemple les langages intrinsèquement parallèles que sont \textsf{Chapel} ou \textsf{Cilk} voire \textsf{Erlang} ne sont pas enseignés dans la filière PRCD de l'Enseirb-Matmeca, probablement car ils ne sont pas suffisamment répandus. Par ailleurs pour les capacités des compilateurs, \textsf{gcc} comprend les directives de la bibliothèque \textsf{OpenMP}, mais ne les trouve pas automatiquement (bien qu'elles soient souvent très simples), en revanche les compilateurs modernes sont souvent capables d'exploiter eux-mêmes la caractère superscalaire d'un processeur.}.

%Dans le cadre du stage, nous nous sommes alors plutôt intéressés aux possibilités d'un Domain Specific Embedded Language (DSEL, équivalent de \emph{langage spécifique embarqué} en français) qui permet une charge de développement moindre, ainsi qu'une appréhension rapide. En effet il s'agit alors d'une bibliothèque fournie au programmeur dans un langage courant préexistant, ne comportant que quelques fonctions reprenant uniquement les paradigmes dont le programmeur a besoin pour résoudre un type de problème spécifique, et se basant sur des bibliothèques performantes préexistantes. 
Nous avons donc décidé de développer une approche basée sur un Domain Specific Embedded Language (DSEL, équivalent de \emph{langage spécifique embarqué} en français) sous la forme d'une bibliothèque, avec comme objectifs de faciliter la prise en main du programmeur et de réduire sa charge de développement, tout en préservant ou en améliorant les performances de ses programmes.
% à recaser plus tard
%Le DSEL se différencie d'un DSL (équivalent de \emph{langage spécifique}) par le fait qu'il est disponible au sein d'un langage général, et permet donc potentiellement au développeur d'effectuer des tâches connexes à la résolution du problème dans un seul et même langage. Il peut ainsi utiliser plusieurs DSEL dans le même code, et donc résoudre des problèmes complexes avec un formalisme adapté à chacun d'entre eux (et par là-même avoir de meilleures performances). Le DSEL est ainsi une solution intermédiaire efficace, qui nous a semblé adaptée à l'objectif du stage qui est la parallélisation automatique d'une catégorie spécifique de programmes. 
Plutôt que de focaliser ce DSEL sur un type spécifique de parallélisme, ou sur un paradigme comme le font de nombreux outils, ce DSEL est focalisé sur une catégorie de problèmes et doit permettre de cacher justement tout ce qui relève du parallélisme.




%% Intro:
%% -> contexte //nécessité de la programmation parallèle
%% -> problématique //dur d'écrire du code parallèle (même pour le compilateur)
%% -> contribution //avoir de l'expressivité, exemple entre QIRAL et SYCL pour exécuter des stencils

%% -*- eval: (flyspell-mode 1); -*-

\chapter{Extraction et formalisation du parallélisme}

Le parallélisme dans un programme informatique est difficile à trouver automatiquement et la plupart du temps il est exprimé explicitement par le développeur. Il est alors nécessaire de lui offrir des outils permettant d'écrire le code parallèle simplement, en lui proposant notamment des fonctions et des mots-clés dans le langage choisi qui sont en accord à la fois avec les besoins matériels et à la fois avec les besoins algorithmiques. Ce chapitre dresse d'abord un état de l'art des principales solutions existantes, puis introduit les principaux objectifs du DSEL développé pendant le stage.

\section{Contexte}

Les applications parallèles sont utilisées depuis un certain nombres d'années afin d'obtenir des performances (en matière de rapidité, et en matière de taille du problème traité). Or il est difficile d'écrire des codes parallèles, ce qui a été évoqué au chapitre précédent. À ce titre divers outils pour faciliter la parallélisation sont apparus, tels que \textsf{ParaScope} dès $1989$ \cite{Art24} qui est éditeur de code interactif. Ces outils se décomposent en de nombreuses catégories à la fois par rapport aux utilisateurs qu'ils visent, et à la fois par rapport au parallélisme qu'ils permettent. Le cas d'application des stencils -- qui est le cas d'application du stage -- est alors explicité afin de voir comment interviennent les différents niveaux de parallélisme. 

\subsection{Expression du parallélisme}

Décrire les problèmes à résoudre et le parallélisme qu'ils peuvent utiliser peut se faire de différentes manières avec les moyens actuels. Notamment il est possible de regrouper quatre catégories qui sont les langages dédiés aux scientifiques (non spécialistes du HPC), les langages et compilateurs pour les développeurs de métier (mais non obligatoirement spécialistes du HPC), l'ensemble des bibliothèques et supports d'exécution spécialisés (destinés aux spécialistes du HPC), et en enfin les outils de formalisation des algorithmes employés (destinés aux spécialistes du HPC).

\paragraph{Langages dédiés aux scientifiques.}
Les langages dédiés aux scientifiques sont la plupart du temps des outils de simulation à destination d'experts dans ce domaine. Dans le cas des mathématiques par exemple, de nombreux logiciels définissant leur propre langage existent : \textsf{\href{http://www.wolfram.com/mathematica/}{Mathematica}}, \textsf{\href{http://fr.mathworks.com/}{MATLAB}} ou encore \textsf{\href{http://pari.math.u-bordeaux.fr/}{PARI/GP}} ; et une partie de leurs algorithmes sont parallélisés (au moins au niveau d'un processeur) pour des raisons de performances. Des initiatives plus poussées du côté du parallélisme sont également présentes en recherche comme \textsf{QIRAL} \cite{Art13,Art14} qui à partir d'algorithmes décrits en \LaTeX permet de générer un code parallèle optimisé pour les processeurs des ordinateurs courants -- mais pas pour les supercalculateurs --, il se base pour cela sur le système de réécriture \textsf{Maude} \cite{Art12}. La logique est alors la suivante : offrir à l'utilisateur les mêmes fonctions qui sont connues dans son domaine, en cachant l'implantation et la parallélisation de ces fonctions (telles qu'un calcul de facteurs premiers par exemple).

\paragraph{Langages dédiés au parallélisme.}
Quant il s'agit de développeurs souhaitant écrire du code parallèle sans forcément connaître les subtilités des machines d'exécutions, différents langages spécifiques ont été créés. De manière générale il s'agit en fait d'un langage déjà connu (\textsf{Java} ou \textsf{C/C++}) auquel sont ajoutés quelques mots-clés ou éléments de syntaxes (souvent pour les tableaux notamment) ; \textsf{\href{http://chapel.cray.com/}{Chapel}}, \textsf{\href{https://www.cilkplus.org/}{Cilk}}, \textsf{\href{http://upc.lbl.gov/}{UPC}}, \textsf{Lime} \cite{Art8}, \textsf{\href{http://x10-lang.org/}{X10}} et le langage proposé dans\cite{Art10} rentrent dans cette catégorie. Alors que les trois premiers cités sont très proches des langages d'origines, les trois derniers sont plus éloignés et proposent des paradigmes différents (comme les Map/Reduce dans \textsf{Lime} ou les horloges dans \textsf{X10}). Avec ces langages viennent souvent les compilateurs associés, toutefois parfois il s'agit plus simplement de réécriture automatique de code (vers un langage déjà existant) comme proposé par \textsf{Par4All} \cite{Ths2}. \textsf{Par4All} peut d'ailleurs être considéré comme un éditeur de code dédié au parallélisme, de même que le projet \textsf{Eclipse \href{http://www.eclipse.org/ptp/}{PTP}} (Parallel Tools Platform)ou les projets présentés dans les articles \cite{Art23,Art25}. Ces éditeurs rassemblent alors généralement plusieurs outils, et sont plus interactifs, mais leur parallélisation n'est pas toujours optimale.

\paragraph{Bibliothèques spécialisées.}
Pour atteindre un parallélisme optimal il est souvent nécessaire de descendre à un niveau plus proche du matériel. Plutôt que d'introduire quelques mots-clés supplémentaires, il s'agit alors de plusieurs dizaines de fonctions regroupées dans des bibliothèques. Il faut alors parfois utiliser différentes bibliothèques afin de gérer les différents niveaux de parallélisme (notamment interprocesseur et intraprocesseur). La bibliothèque utilisée dans le cadre du stage en fait partie : \textsf{\href{https://www.khronos.org/sycl}{SYCL}}. C'est celle qui est la plus haut niveau tout en offrant le plus de possibilités. La bibliothèque connue qui s'en rapproche le plus est \textsf{\href{https://www.khronos.org/opencl/}{OpenCL}} mais nécessite déjà plus d'efforts notamment pour gérer les communications entre les processeurs. En revanche \textsf{\href{http://openmp.org/wp/}{OpenMP}} (qui est en fait intégrée dans la plupart des compilateurs pour \textsf{Fortran}, \textsf{C/C++}) est beaucoup plus simple mais ne s'utilise que pour les processeurs génériques et non les processeurs graphiques. D'autres bibliothèques sont justement dédiés à la gestion de plusieurs types de processeurs en même temps, il s'agit des \emph{supports d'exécutions} tels que \textsf{\href{http://starpu.gforge.inria.fr/}{StarPU}} ou \textsf{\href{http://icl.utk.edu/parsec/}{PaRSEC}}. Comme la plupart des problèmes parallélisés sont liés à l'utilisation de grand tableaux, \textsf{Kokkos} \cite{Web5} et \textsf{\href{http://hpc.pnl.gov/globalarrays/}{Global Arrays}} permettent de simplifier leur utilisation et répartition sur les différents processeurs, mais ils ne gèrent pas le lancement des calculs contrairement aux supports d'exécutions. Plus spécifiquement encore l'article \cite{Art2} propose une solution simplifiant l'utilisation des tableaux avec le parallélisme, au sein d'un processeur. Enfin certaines rares bibliothèques se penchent plutôt sur les caractéristiques des algorithmes utilisés comme \textsf{SkePU} \cite{MstThs1}, qui s'apparente à un moteur d'exécution et utilise la formalisation de squelettes algorithmiques.


\paragraph{Formalisation des algorithmes.}
La formalisation des algorithmes parallèles se fait à travers différents outils comme la représentation polyédrique, utilisée lorsqu'il s'agit de parcours de tableaux au sein de boucles affines. Cette représentation permet d'ailleurs de mieux classifier les algorithmes, notamment en squelettes comme dans \textsf{SkePU} et \textsf{Bones} \cite{Art4,Art3,Art9}. Ces squelettes permettent de définir le motif d'accès aux données des tableaux impliqués dans le calcul de façon précise et ainsi de l'optimiser. Un tel exemple de motif -- plutôt général -- est le paradigme de \emph{Map/Reduce} qui est proposé par \textsf{SkePU} (ainsi que par le langage dédié \textsf{Lime} cité plus-haut). Cependant tous les algorithmes ne décrivent pas un calcul : certains sont au contraire destinés à ordonnancer les calculs, c'est-à-dire les répartir sur les différents processeurs disponibles. Encore plutôt à l'état de concept, c'est dans ce but que sont utilisés les \emph{composants} permettant les prédictions (et optimisations) du temps d'exécution tels que dans l'article \cite{Art7}, ainsi que les \emph{conteneurs intelligents} pour la gestion des communications avec la mémoire (et aussi les prédictions de temps) dans \textsf{SkePU} \cite[p.~78]{Ths1}.

\subsection{Cas d'application: les stencils}
\label{sec:stencil_base}

De nombreux problèmes notamment physiques ou mathématiques utilisent un certain type de calculs sur des grilles d'éléments qui sont regroupés sous le nom de \emph{stencils}. Dans ces problèmes chaque élément au sein d'une grille est calculé en fonction des valeurs de ses voisins. En une dimension, il s'agirait par exemple d'un tableau dans lequel chaque élément reçoit la somme de son voisin de droite et de celui de gauche. La notion de stencil est générique pour toute dimension, pour tout nombre de voisins impliqués, ainsi que pour leur disposition. La figure \ref{fig:stencil_base} présente un exemple de stencil en deux dimensions où seuls les quatre voisins immédiats interviennent dans le calcul de chaque élément bleu de la grille. Comme cette forme de stencil nécessite l'utilisation de points dans les deux sens de chaque dimension, une partie des points ne peuvent pas être mis à jour. Cette zone est appelée la zone \emph{fantôme} puisque qu'aucun calcul n'est fait sur ses éléments mais est obligatoire car ses éléments sont utilisés pour le calcul des éléments internes à la grille. Cette zone fantôme est représentée avec des éléments rouges sur la figure \ref{fig:stencil_base}.

\begin{figure}[!bp]
\floatbox[{\capbeside\thisfloatsetup{capbesideposition={right,center},capbesidewidth=4cm}}]{figure}[\FBwidth]
{\caption{Exemple d'une grille en deux dimensions avec un stencil à 4 points.}\label{fig:stencil_base}}
{\begin{tikzpicture}
\tikzstyle{place}=[circle,thick,draw=blue!75,fill=blue!20,minimum size=2mm]

\node[place] (p1) at (0,0) {};
\node[place] (p2) at (0,1) {};
\node[place] (p3) at (1,0) {};
\node[place] (p4) at (-1,0) {};
\node[place] (p5) at (0,-1) {};

\node[place] () at (-1,-1) {};
\node[place] () at (-1,1) {};
\node[place] () at (1,-1) {};
\node[place] () at (1,1) {};

\foreach \x in {-2,...,2}
{
  \node[place] () at (\x,2) {};
  \node[place] () at (\x,-2) {};
  \node[place] () at (2,\x) {};
  \node[place] () at (-2,\x) {};
}

\draw[->] (p2) to (p1);
\draw[->] (p3) to (p1);
\draw[->] (p4) to (p1);
\draw[->] (p5) to (p1);

\tikzstyle{place}=[circle,thick,draw=red!75,fill=red!20,minimum size=2mm]

\foreach \x in {-3,...,3}
{
  \node[place] () at (\x,3) {};
  \node[place] () at (\x,-3) {};
  \node[place] () at (3,\x) {};
  \node[place] () at (-3,\x) {};
}



\end{tikzpicture}
}
\end{figure}

L'importante utilisation des stencils en recherche a conduit à des DSEL comme \textsf{Pochoir} \cite{Art18} ou  \textsf{Blitz++} \cite{Art5}, et des langages tels que \cite{Art19}, qui -- hormis \textsf{Blitz++} nécessitent leur propre compilateur afin d'effectuer des optimisations. Toutefois les stencils sont aussi utilisées dans des applications physiques non scientifiques, par exemple dans le calcul des graphismes de jeux vidéos \cite{Art15}. Dans les deux cas la performance est primordiale et il est donc nécessaire d'exploiter les possibilités de parallélisme des machines utilisées pour ces calculs.  

Si l'on suppose que la mise à jour de la grille par application du stencil n'est pas calculée en place -- c'est-à-dire que l'on dispose d'un tableau d'entrée dans lequel les éléments sont seulement lus, et d'un tableau de sortie où les résultats du calcul du stencil sur chaque élément sont stockés -- le stencil peut être appliqué à chaque élément bleu en même temps. Le parallélisme se trouve donc de manière très aisée dans ce type de problème. Avec une  machine possédant autant de cœurs que d'éléments dans la grille, la mise à jour de la grille se fait donc en une seule étape et tous les processeurs exécutent les mêmes instructions, simplement sur des éléments différents ; le calcul est donc bien adapté à une carte graphique qui regroupe justement ces caractéristiques : nombre élevé de cœurs et un même code s'exécutant sur tous ceux-ci. Puisque le calcul n'est pas en place, pour l'appliquer plusieurs fois d'affiliés il suffit d'intervertir les grilles d'entrées et de sortir entre chaque mise à jour, ou bien de recopier la sortie vers l'entrée, ce qui est toutefois plus lent.



Parler de la mémoire !


\section{Conception d'un DSEL pour les stencils}

A relire \cite{Art22}

\subsection{Objectifs sémantiques}

De quoi a-t-on besoin pour décrire un stencil

-> indices coefs

-> pattern du stencil

-> tableaux coefs

-> tableaux données

-> abstraire les structurations de données (copies tuiles, et data layout)

\subsection{Objectifs de performances}


-> tuiles/caches (data reuse)

-> calculs à la compilation


QCD \cite{Art1}
3.5D blocking \cite{Art11}


%% Chap2:
%% -> contexte
%% -> état de l'art
%% -> problématique

%% Chap3:
%% -> proposition (aspect théorique)

%% -*- eval: (flyspell-mode 1); -*-

\chapter{Bibliothèque de stencils basée sur \textsf{SYCL}}

Au cours du stage a été développée un DSEL de stencils sous la forme d'une bibliothèque \textsf{C++} miniature. Ce DSEL essaie de reprendre les objectifs explicités dans le chapitre précédent, toutefois il s'agit plutôt d'une preuve de faisabilité ce pour quoi tous les objectifs n'ont pas été implémentés (mais leur intégration a été prévue). Ce chapitre se concentre sur la présentation de ce DSEL en ce qui concerne son utilisation, quelques détails d'implantations, ainsi que le travail restant.

\section{Présentation générale}

Le DSEL implémenté s'appuie sur très peu de classes afin de simplifier son utilisation, il s'agit essentiellement de deux classes pour les coefficients (constants ou variables), plusieurs classes pour la description d'une opération (tableaux d'entrées et sorties, fonctions d'accès aux éléments ou aux coefficients telles que décrites dans la section \ref{sec:obj_sem}) et une unique classe pour lancer les calculs. Concrètement l'écriture d'un problème à base de stencils dans le DSEL se fait en cinq étapes :

\begin{enumerate}
\item initialiser les coefficients ;
\item les assembler pour former un stencil au sein d'un objet \emph{stencil};
\item regrouper les informations nécessaires au calcul (stencil, tableaux et fonctions) au sein d'un objet \emph{operation};
\item déclarer une queue de calcul dans le formalisme \textsf{SYCL} ;
\item appeler la méthode de lancement des calcul sur l'objet \emph{operation}.
\end{enumerate}

\subsection{Création des stencils}

Les coefficients peuvent être créés en version constante ou variable, avec les indices de décalage spécifiés dans le \emph{template} de la classe, ils seront donc statiques. La logique de construction n'est pas exactement la même pour les deux classes puisque la version constante nécessite de préciser le type du coefficient (entier, flottant, etc\ldots). Enfin notons que seule la version en deux dimensions a été implémentée et est ici présentée.
\begin{listing}[H]
\caption{Création de coefficients, avec leur indice de décalage.}
\begin{minted}[frame=single]{cpp}
/* la valeur du coefficient est affectée à la construction */
coef_fxd2D<-1, 1, double> c2(0.2);
/* les valeurs du coefficient nécessiteront une fonction tierce */
coef_var2D<0, 0> c1;
\end{minted}
\end{listing}

Une fois les coefficients initialisés, il est possible de les rassembler pour former le motif du stencil. Il est d'ailleurs possible de combiner plusieurs stencils ensembles de la même manière. Toutefois il n'est possible d'agréger que des coefficients constants et du même type entre eux, ou bien que des coefficients variables entre eux, il n'est pas possible de mélanger les deux possibilités.
\begin{listing}[H]
\caption{Agrégation des coefficients pour former un stencil à cinq points.}
\begin{minted}[frame=single]{cpp}
/* création d'un stencil à 5 points */
auto st = c1+c2+c3+c4+c5; 
\end{minted}
\end{listing}
Il est important de remarquer deux choses sur cet extrait de code : l'usage du mot-clé \textsf{auto} et l'usage de l'opérateur \textsf{+}. Le mot-clé \textsf{auto} permet de cacher à l'utilisateur le vrai type du stencil qui, étant un template, contient de nombreuses informations et serait donc fastidieux à écrire. Concrètement écrire directement le type du stencil reviendrait à écrire les propriétés du stencil sous une forme très peu lisible (tous les indices étant les uns à la suite de autres sur la même ligne), ce qu'il vaut donc mieux éviter. Quant au \textsf{+} il est utilisé pour faire la \og concaténation \fg~des coefficients mais n'a rien à voir avec l'opération mathématique qui est effectuée derrière. Lier les deux pourraient être intéressant, mais complexe à mettre en place car dans la plupart des cas, les opérations utilisées sont des additions.

Les coefficients variables nécessitent l'initialisation d'un tableau tiers, dont le type est basé sur \textsf{SYCL}. Il s'agit d'un \emph{buffer} à une seule dimensions comprenant n'importe quel nombre d'éléments. Ce buffer est alors associé à une fonction d'accès, qui doit veiller à ne pas dépasser les limites du tableau (non vérifié par la compilation). L'exemple de code suivant présente l'équivalent d'un coefficient constant sous forme de coefficient variable.
\begin{listing}[H]
\caption{Exemples de fonction d'accès aux coefficients.}
\begin{minted}[frame=single]{cpp}
/* valeur du coefficient */
float tab_var = 1.0;
/* création du buffer */
coefBuffer = sycl::buffer<float,1>(&tab_var, sycl::range<1> {1});
/* création de la fonction d'accès */
float  fac_coef(int a,int b, int c, int d, 
                sycl::accessor<float, 1, sycl::access::read>  acc) {
  return acc[0];
}
\end{minted}
\end{listing}
Bien sûr dans le cas d'un coefficient variable, la fonction d'accès peut être beaucoup plus complexe et peut notamment faire intervenir les quatre variables \verb!a, b, c, d! qui sont les indices globaux puis de décalage de l'élément en cours de calcul, tandis que \verb!acc! est le tableau des coefficients. L'écriture de ces fonctions d'accès n'est pas vraiment aisé, mais elles sont nécessaires dans certains types de calculs comme la chromodynamique quantique.

\subsection{Application des stencils}

Avant de déclencher le calcul, il reste à préciser les informations de l'opération : tableaux d'entrée et sortie notamment, ainsi que leur fonction d'accès. Ces fonctions d'accès sont plus simples puisqu'elles ne font intervenir que les indices globaux, en revanche il est nécessaire d'en définir une pour la sortie et une pour l'entrée, même si les deux accèdent aux éléments de la même façon. Ces deux versions différents sont nécessaires car la fonction du tableau de sortie, dans lequel sont écrits les résultats, doit renvoyer des éléments modifiables et non seulement en lecture seule. Deux prototypes sont donnés en exemple ci-dessous.
\begin{listing}[H]
\caption{Exemples de fonctions d'accès aux éléments.}
\begin{minted}[frame=single]{cpp}
/* Fonction d'accès pour la sortie, notez le & */
float& fdl_out(int a,int b, 
               sycl::accessor<float, 2, sycl::access::write> acc) 
               {return acc[a][b];}
/* Fonction d'accès pour l'entrée */
float  fdl_in(int a,int b, 
              sycl::accessor<float, 2, sycl::access::read>  acc) 
              {return acc[a][b];}
\end{minted}
\end{listing}

Les tableaux de l'entrée et de la sortie sont déclarés sous forme de buffers, exactement comme dans le cas précédent à la différence qu'ils sont bien en deux dimensions cette fois. Une fois tous les éléments déclarés, il est donc possible de former l'opération comme dans le code ci-dessous.
\begin{listing}[H]
\caption{Assemblage des informations d'entrées et sorties pour former une opération.}
\label{lst:tab_io}
\begin{minted}[frame=single]{cpp}
/* assemblage des éléments de l'entrée, dans l'ordre :
 * type, buffer des éléments en entrée, buffer des 
 * coefficients, fonction d'accès aux éléments,
 * fonction d'accès aux coefficients
 */
input_var2D<float, &ioABuffer, &ioBBuffer, &fdl_in, &fac> work_in;
/* assemblage des éléments de sortie, dans l'ordre :
 * type, buffer des éléments en sortie, fonction
 * d'accès aux éléments
 */
output_2D<float, &ioBuffer, &fdl_out> work_out;
/* création de l'opération */
auto op_work = work_out << st << work_in;
\end{minted}
\end{listing}
Le code pour les coefficients constants est similaire, seule la description de l'entrée est légèrement simplifiée puisqu'il n'y pas besoin de la fonction d'accès aux coefficients ni du buffer associé. 

Enfin il reste à appeler le calcul, c'est l'étape la plus simple qui n'appuie que sur deux méthodes différentes, adaptées suivant les besoins d'optimisations de l'utilisateur (avec tuilage ou non). Pour lancer les calculs il est nécessaire de déclarer une queue qui contient les informations relatives au matériel disponible et est gérée par \textsf{SYCL}. La découverte du matériel est automatique, la simple déclaration de la variable permet son utilisation. Nous obtenons alors le code de lancement suivant.
\begin{listing}[H]
\caption{Lancement d'opérations de calculs.}
\begin{minted}[frame=single]{cpp}
{   
  /* création de la queue */
  sycl::queue myQueue; 
  /* application d'une opération avec tuilage */
  op_work.doLocalComputation(myQueue);
  /* application d'une autre opération sans tuilage */
  op_copy.doComputation(myQueue);
}
\end{minted}
\end{listing}

Dans tout le code présenté, la taille de la grille n'apparaît qu'une seule fois : lors de la création des buffers. Le DSEL se charge de déterminer les zones fantômes et applique automatiquement le stencil sur le plus d'éléments possibles (les tableaux d'entrées et sorties doivent avoir les mêmes dimensions). Ceci est un bénéfice pour l'utilisateur qui n'a pas besoin de répéter des informations redondantes, mais aussi un risque pour lui car il est alors susceptible d'oublier l'existence même de ces zones fantômes. 

\section{Détails d'implantation}

L'utilisation du \textsf{C++} et notamment de la méta-programmation a permis le résultat présenté ci-dessus. Le fonctionnement interne du DSEL est volontairement caché à l'utilisateur, nous le détaillons donc quelque peu dans cette section, notamment pour la gestion des variables statiques qui est l'un des intérêts principaux de la méta-programmation par template. Cela permet alors de spécialiser les fonctions lors de la compilation et donc de les rendre plus rapides lors de l'exécution. Concrètement la méta-programmation permet d'écrire un code qui sera modifié (ou plus exactement spécialisé) par d'autres parties du code. Par exemple en \textsf{C++} une classe template sera spécialisée pour chacune de ses instances déclarées dans le code et ayant des paramètres différents. La méta-programmation est au cœur de tout le développement du DSEL, elle permet de se placer \og au-dessus \fg~du langage à partir du langage lui-même. Notons que ce principe se retrouve dans d'autres langages informatiques tel que \textsf{Maude}, et même dans les langages naturels \cite{Web6}.

\subsection{Précisions sur le \textsf{C++14} et l'utilisation de \textsf{SYCL}}

Le DSEL est développé avec le dernier standard officiel du \textsf{C++}, \textsf{C++14}, afin d'être compatible avec \textsf{SYCL}. En effet l'implantation utilisée de \textsf{SYCL} est \textsf{triSYCL}, qui est une bibliothèque écrite en \textsf{C++14} par Ronan Keryell. Celle-ci utilise notamment les possibilités des \emph{lambda fonctions} et du mot-clé \textsf{auto}, introduits à partir de \textsf{C++11} et augmentés en \textsf{C++14}. Outre les éléments de langages, \textsf{triSYCL} s'appuie également sur \textsf{boost} par exemple pour la gestion des tableaux multidimensionnels, ainsi que sur \textsf{OpenMP} pour une parallélisation naïve. L'interface \textsf{SYCL} fournit un certain nombre de fonctions pour faciliter le parallélisme, parmi elles notons :
\begin{itemize} 
\item la découverte (potentiellement) automatique du matériel disponible sur la machine ;
\item la gestion automatique des transferts de données avec ces matériels (la mémoire globale des cartes graphiques principalement) ;
\item le lancement de \emph{noyaux} de calcul sur le matériel désiré, avec configuration possible des tailles de \emph{warps}.
\end{itemize}
L'interface contient de nombreuses autres possibilités, pour beaucoup calquées sur \textsf{OpenCL}, mais seule une petite partie a pu être testée puis utilisée dans le DSEL, car l'implantation de \textsf{triSYCL} n'est pas terminée. Toutefois l'implantation actuelle est suffisante pour les besoins du DSEL.

En ce qui concerne le DSEL lui-même, il utilise la notion d'\emph{expression template} \cite{Web4,Art21} qui est une façon de décrire des expressions grâce à la méta-programmation. C'est ainsi que sont stockés les stencils : il s'agit d'un arbre binaire de coefficients. Les feuilles de l'arbre sont les coefficients tandis que les nœuds sont des sous-expressions de la racine (donc des morceaux du motif du stencil). Dans le cas du code \ref{lst:st_4pts} ci-dessous présentant l'agrégation de quatre coefficients, l'arbre construit est en forme de peigne (à droite ou à gauche), comme dans la figure \ref{fig:tree1_coef}.
\begin{listing}[H]
\caption{Code de création d'un stencil à 4 points.}
\label{lst:st_4pts}
\begin{minted}[frame=single]{cpp}
/* création d'un stencil à 4 points */
auto st = c1+c2+c3+c4; 
\end{minted}
\end{listing}
\begin{figure}[!h]
\floatbox[{\capbeside\thisfloatsetup{capbesideposition={right,center},capbesidewidth=4cm}}]{figure}[\FBwidth]
{\caption{Exemple d'arbres en peigne d'un stencil à $4$ points formé en une seule fois, par le code \ref{lst:st_4pts}.}\label{fig:tree1_coef}}
{\begin{tikzpicture}[scale=0.8]
\begin{scope}
\node (a) {$\oplus$}
  child {node {$\oplus$}
    child {node {$\oplus$}
      child {node {$\oplus$}
        child {node {$c_1$}}
      }
      child {node {$c_2$}}
    }
    child {node {$c_3$}}
  } 
  child {node {$c_4$}};

\end{scope}

\begin{scope}[xshift=6cm]
\node (b) {$\oplus$}
  child {node {$c_1$}}
  child {node {$\oplus$}
    child {node {$c_2$}}
    child {node {$\oplus$}
      child {node {$c_3$}}
      child {node {$\oplus$}
        child {node {$c_4$}}
      }
    }
  } ;

\end{scope}

\path (a) -- (b) node [midway] {ou};

\end{tikzpicture}
}
\end{figure}
Si le stencil est formé en deux temps, il est alors possible d'avoir un arbre équilibré comme dans la figure \ref{fig:tree2_coef}, grâce au code \ref{lst:st_2_2pts}.
\begin{listing}[H]
\caption{Code de création d'un stencil par le biais de deux autres stencils.}
\label{lst:st_2_2pts}
\begin{minted}[frame=single]{cpp}
/* création d'un stencil à 4 points */
auto st1 = c1+c2;
auto st2 = c3+c4;
auto st = st1+st2;
\end{minted}
\end{listing}
\begin{figure}[!h]
\floatbox[{\capbeside\thisfloatsetup{capbesideposition={right,center},capbesidewidth=4cm}}]{figure}[\FBwidth]
{\caption{Exemple d'arbre équilibré d'un stencil à $4$ points formé en plusieurs temps, par le code \ref{lst:st_2_2pts}.}\label{fig:tree2_coef}}
{\begin{tikzpicture}[scale=0.8,level/.style={sibling distance=40mm/#1}]
\node (a) {$\oplus$}
  child {node {$\oplus$}
    child {node {$c_1$}}
    child {node {$c_2$}}
  }
  child {node {$\oplus$}
    child {node {$c_3$}}
    child {node {$c_4$}}
  };

\end{tikzpicture}
}
\end{figure}
Notons que la forme de ces arbre n'a pas été étudiée pour l'obtention des performances, mais il est tout à fait possible que cette forme les influencent car elle influe sur l'ordre de calcul des opérations dans le stencil. 

En supposant que certaines formes sont plus performantes que d'autres, il est primordial de pouvoir alors spécialiser le code en ayant en interne à la bibliothèque des fonctions optimisées pour certains types de motifs. Cela est rendu possible en \textsf{C++} grâce à la possibilité de redéfinir des classes et méthodes plusieurs fois avec des paramètres de template différents : la spécialisation partielle \cite{Web1}. En pratique il est même possible d'avoir des template à l'intérieur de template, ce qui permet un déclenchement retardé de la spécialisation d'une classe par rapport à ses méthodes. Cette méthode est notamment utilisée pour éviter à l'utilisateur de préciser le type dans les coefficients variables, celui-ci est propagé à la méthode lors de la construction des tableaux de coefficients (cf l'extrait de code \ref{lst:tab_io}). La spécialisation partielle \emph{retardée} est relativement complexe, même pour le compilateur, d'où l'introduction d'une syntaxe peu compréhensible \cite{Web2}. En revanche l'utilisateur doit préciser le type dans le cas des coefficients constants. Cela est redondant pour lui -- il passe l'information deux fois --, mais aussi pour le DSEL car tout le code doit alors être dupliqué, et possiblement de manière non statique ! En effet il est alors nécessaire de stocker la valeur du coefficient dans l'objet-même donc tous les autres objets (tels qu'une opération) l'utilisant doivent avoir un lien \emph{a-un} vers lui, qui est a priori dynamique. En pratique le compilateur est parfois capable de faire une telle optimisation, mais cela n'est pas garanti, et n'a pas été vérifié.


\subsection{Calcul statique des paramètres de tuilage}
\label{sec:param_tuile}

Plusieurs paramètres configurent le tuilage, notamment la taille de la zone fantôme ainsi que la taille maximum d'un warp. Ces paramètres sont calculés lors de la compilation, toujours grâce à la méta-programmation. Pour la taille de la zone fantôme, le calcul est très simple, il est fait lors de la construction du stencil ou plus exactement lors de l'agrégation des coefficients (ou stencils) entre eux. En effet les objets coefficients et stencils contiennent tous les deux des indices locaux étant le maximum de décalage dans chacune des directions (quatre en deux dimensions). Le maximum est mis-à-jour à chaque agrégation, en fonction des maximums des fils dans l'arbre binaire du stencil. C'est grâce à ce maximum (un dans chaque direction) qu'est déterminée la zone fantôme qui est définie par l'\emph{offset} à partir duquel la zone de calculs commence, ainsi que la largeur totale de la zone fantôme dans chaque dimension. En deux dimension, le stencil à quatre points est donc large de $2$ éléments dans chacune des deux dimensions, tandis que l'offset est de $1$ dans chacune des deux dimensions également.

L'ensemble de ce calcul est fait de manière statique, car ces paramètre sont liés au template des classes impliquées qui elles-même sont statiques -- dans le cas des coefficients variables. Ce fonctionnement a été validé par la décompression d'un fichier binaire censé affiché le résultat de ce calcul, on peut y vérifier que le résultat y apparaît directement dans le code, sans calcul le précédent. Ce résultat peut sembler paradoxal : c'est avec les coefficients variables que nous obtenons des paramètres statiques. Cela n'est pas contradictoire car ce qu'il se produit dans les cas des coefficients variables est une séparation totale entre l'objet définissant uniquement les indices de décalage, et l'objet permettant d'accéder aux coefficients. Par simplicité pour l'utilisateur, ces deux objets sont réunis en un seul dans le cas d'un coefficient fixe, mais à un certain prix : celui de l'optimisation statique. Il s'agit là d'un parfait exemple de l'équilibre à trouver entre les trois piliers de Todd Veldhuizen pour une librairie optimisée : beauté, sûreté et rapidité.

Une fois la zone fantôme calculée, et ainsi le nombre d'éléments à charger que va nécessiter le calcul d'une macro-tuile, il est alors possible de déterminer la taille optimale de la macro-tuile. Cette taille est déterminer de manière récursive lors de la compilation en fonction de trois paramètres configurés dans les headers, et qui sont en pratique des variables d'environnement liées au matériel : nombre de cœurs maximum dans chaque dimension pour un warp de la carte graphique, quantité maximale de mémoire locale disponible pour un warp. Il s'agit alors d'obtenir la plus grande taille possible de macro-tuile, sans dépasser la quantité de mémoire disponible. Le principe de l'algorithme est alors très simple : en commençant par la taille minimale (celle qu'occupe la zone fantôme) l'algorithme essaie d'incrémenter le nombre d'éléments dans une dimension (ou les deux si possible) jusqu'à ce que l'une des conditions suivantes soit fausse : soit le nombre de cœurs n'est plus suffisant, soit il ne reste plus assez de mémoire. Du point de vu de l'implantation en \textsf{C++} il s'agit d'une classe template qui possède un lien \emph{a-un} avec elle-même, à la différence près que celle qu'elle possède stocke une taille de macro-tuile incrémentée de un (dans une ou deux dimensions) ainsi qu'un booléen pour signifier l'arrêt ou non de la récursion. Suivant cette méthode, la macro-tuile ne s'adapte pas à la forme du stencil, mais cela est toujours possible grâce à la spécialisation partielle évoquée dans la sous-section précédente.

De même que pour la zone fantôme, le tuilage est effectué de manière statique et le résultat du calcul est bien présent directement dans le binaire exécutable du programme. Les calculs statiques ne concernent d'ailleurs pas que des paramètres, mais aussi la sémantique du DSEL. En effet la présence de \textsf{static\_assert} dans le code permet d'éviter à l'utilisateur d'écrire du code fallacieux, en déclenchant une erreur lors de la compilation. Ainsi pour le tuilage, si le stencil est trop large pour la taille d'une macro-tuile, une erreur va se déclencher. De même l'agrégation de coefficients variables avec des coefficients fixes va déclencher une erreur lors de la compilation. Cela est très utile pour le développeur car il est beaucoup plus simple de corriger un problème lorsqu'il est détecté à la compilation que lorsqu'il est détecté à l'exécution, c'est de plus un gain de temps dans le développement. 

\section{\'Evolutions possibles}
\label{sec:evol_biblio}

Le DSEL développé est une preuve de concept, ainsi nous savons maintenant comment développer certains aspects afin de rendre la librairie compétitive avec d'autres librairies du point de vue des performances, et peut-être même meilleure du point de vue de l'expressivité du langage -- le caractère joli. Parmi les fonctionnalités suggérées, certaines se situent au niveau de l'utilisateur et des possibilités offertes -- le \emph{front-end} -- , d'autres sont cachées et concernent les possibilités de parallélisme et d'optimisations -- le \emph{back-end}.

\subsection{Améliorations du front-end}

Le front-end doit offrir la possibilité à l'utilisateur de décrire la plupart des problèmes utilisant des stencils. À ce titre il est nécessaire de rajouter la possibilité de spécialiser la fonction d'application d'un coefficient sur un élément, représentée par le symbole $\otimes$ défini à la section \ref{sec:obj_sem}. En effet cela permettrait non seulement de gérer des types plus complexes que les entiers ou les flottants (par exemple des matrices ou des complexes), mais aussi de traiter des problèmes non linéaires (par exemple la mise au carré des éléments d'une grille, se faisant par un stencil à un seul point donc). Le rajout d'une telle fonction est très simple, il suffit de la mettre en paramètre du template du coefficient voulu (à la suite des indices de décalage). De même il serait envisageable de pouvoir modifier l'opérateur de réduction $\oplus$, mais ce cas est moins utile et moins intuitif de manière sémantique. En effet le plus simple serait de la rajouter en paramètre du template contenant les tableaux d'entrées du stencil et fonctions associées, mais cet objet deviendrait alors un objet \og fourre-tout \fg~ sans véritable sens. 

Par ailleurs les calculs de stencils sont souvent chaînés les uns aux autres : le tableau de sortie de l'application d'un stencil est le tableau d'entrée du stencil suivant. Ces traitements successifs ne devraient pas obliger l'utilisateur à redéfinir des nouveaux tableaux d'entrées et sorties pour chaque stencil. Il serait plus simple de pouvoir créer une unique opération pour toute la chaîne de stencils, comme dans le code \ref{lst:agreg1_stencil} suivant. Du point de vue l'implantation et si le tuilage est utilisé, il est alors nécessaire de prévoir des communications de zones fantômes entre les warp entre l'application de chaque stencil ou alors d'utiliser la technique présentée dans la section suivante.

\begin{listing}[H]
\caption{Exemple de syntaxe pour le chaînage de plusieurs stencils sur une même entrée.}
\label{lst:agreg1_stencil}
\begin{minted}[frame=single]{cpp}
/* assemblage des éléments de l'entrée, dans l'ordre */
input_var2D<float, &ioABuffer, &ioBBuffer, &fdl_in, &fac> work_in;
/* assemblage des éléments de sortie */
output_2D<float, &ioBuffer, &fdl_out> work_out;
/* création de l'opération */
auto op_work = work_out << st3 << st2 << st1 << work_in;
\end{minted}
\end{listing}

Un autre type d'agrégation est également envisagé, il s'agit d'appliquer des stencils différents (vers des sorties différents donc) sur une même grille d'entrée. Le gain en performance est alors potentiellement élevé car le taux de réutilisation des données est ainsi augmenté : le tableau d'entrée est chargé une seule fois et sert à plusieurs stencils en même temps. Une syntaxe possible est présentée dans le code \ref{lst:agreg2_stencil}, il s'agit alors de placer plusieurs objets d'éléments de sorties les uns à la suite des autres, entrelacés avec les stencils concernés. 

\begin{listing}[H]
\caption{Exemple de syntaxe pour l'application sur une même entrée de plusieurs stencils différents vers leurs sorties respectives.}
\label{lst:agreg2_stencil}
\begin{minted}[frame=single]{cpp}
/* création de l'opération */
auto op_work = work_out3 << st3 << work_out2 << st2 << work_out1 << st1 << work_in;
\end{minted}
\end{listing}

Une telle sémantique nécessite la vérification poussée des tableaux en jeux, afin de détecter lors de la compilation que l'utilisateur n'a pas utilisé plusieurs fois le même tableau de sortie, ou a utilisé un même tableau en sortie et en entrée. Comme le nombre de stencils que l'utilisateur peut vouloir agréger est inconnu à l'avance, il n'est pas possible d'utiliser la technique de spécialisation partielle directement (car il faudrait alors écrire dans le code du DSEL tous les cas d'erreurs possibles, ce qui est du coup impossible). Il convient alors d'utiliser les \emph{variadic template} introduit au \textsf{C++11} \cite{Art6}, qui permettent de créer des listes statiques lors de la compilation, et par exemple d'en compter les éléments, mais aussi de pouvoir les différencier. En pratique cette technique est difficile à mettre en œuvre et nécessite beaucoup de calculs tandis que la limite de pile de récursion des template lors de la compilation se situe autour de $900$ appels. Toutefois il est peu probable que l'utilisateur utilise ne serait-ce que $10$ sorties différentes, ce qui même avec un algorithme naïf quadratique (comparaison de la liste avec elle-même) ne provoquerait que $100$ appels récursifs.


\subsection{Améliorations du back-end}
\label{sec:amelior_backend}

Dans le cas du calcul classique de Jacobi en deux dimensions (qui peut servir notamment à flouter des images), le même stencil est appliqué un certain nombre de fois d'affilée. Or cela nécessite beaucoup de communications, qu'il est possible de diminuer en effectuant en contrepartie quelques calculs supplémentaires redondants (et en utilisant un peu plus de mémoire). Selon les coûts de communications, les coûts de calculs et la mémoire disponible, cette technique peut fournir un gain substantiel de performances, il s'agit de l'\emph{overlapped tiling} ou tuilage chevauchant. Cette technique est assez développée, et est même encore adaptée à la hiérarchie actuelle des ordinateurs \cite{Art17}. Plus concrètement, ce principe est simplement une extension du tuilage tel que représenté dans la figure \ref{fig:stencil_mtile} au chapitre précédent. Les macro-tuiles se chevauchent puisqu'elles nécessitent une zone fantôme qui se trouve chez leurs voisines. Cette zone est prévue pour une seule application dans le cas de la figure \ref{fig:stencil_mtile}, mais il peut très bien il y en avoir plusieurs, concentriques, dont les éléments calculés sur les bords seront redondants. La figure \ref{fig:stencil_mtile2} représente les calculs effectifs pour deux mises-à-jour consécutive (en place, sans communications ni synchronisation) d'une macro-tuile. La macro-tuile centrale contient les résultats tandis que le premier contour contient les résultats (en vert) de la première application du stencil, utilisés comme fantômes lors de l'application du second stencil. Le contour en pointillés le plus éloigné contient les éléments fantômes (en rouge) de l'application du premier stencil. Les données utilisées par la macro-tuile débordent alors largement sur les macro-tuiles voisines dans le cas d'une si petite taille de macro-tuile. Il convient en pratique d'utiliser le tuilage chevauchant pour l'application de quelques stencils seulement (moins d'une dizaine) pour des tuiles de l'ordre de la centaine d'éléments dans chaque dimension.

\begin{figure}[!h]
\floatbox[{\capbeside\thisfloatsetup{capbesideposition={right,center},capbesidewidth=4cm}}]{figure}[\FBwidth]
{\caption{Application d'un stencil en place deux fois de suite sur la macro-tuile centrale. En rouge les éléments fantômes du premier stencil, en vert ceux du second.}\label{fig:stencil_mtile2}}
{\begin{tikzpicture}[scale=0.8]
\tikzstyle{place1}=[circle,thick,draw=blue!75,fill=blue!20,minimum size=4mm]
\tikzstyle{place2}=[circle,thick,draw=green!75,fill=green!20,minimum size=4mm]
\tikzstyle{place3}=[circle,thick,draw=red!75,fill=red!20,minimum size=4mm]

\draw[dashed,fill=orange!20] (2.5,7.5) rectangle (6.5,8.5);
\draw[dashed,fill=orange!20] (2.5,0.5) rectangle (6.5,1.5);

\draw[dashed,fill=orange!20] (1.5,6.5) rectangle (0.5,2.5);
\draw[dashed,fill=orange!20] (8.5,6.5) rectangle (7.5,2.5);

\draw[dashed,fill=orange!20] (1.5,1.5) rectangle (7.5,7.5);


\draw[dashed,fill=orange!40] (2.5,6.5) rectangle (6.5,7.5);
\draw[dashed,fill=orange!40] (2.5,1.5) rectangle (6.5,2.5);

\draw[dashed,fill=orange!40] (2.5,6.5) rectangle (1.5,2.5);
\draw[dashed,fill=orange!40] (7.5,6.5) rectangle (6.5,2.5);


\draw[fill=orange!60] (2.5,2.5) rectangle (6.5,6.5);


\draw (2.5,2.5) -- (-0.5,2.5);
\draw (6.5,2.5) -- (9.5,2.5);

\draw (2.5,6.5) -- (-0.5,6.5);
\draw (6.5,6.5) -- (9.5,6.5);


\draw (2.5,2.5) -- (2.5,-0.5);
\draw (2.5,6.5) -- (2.5,9.5);

\draw (6.5,2.5) -- (6.5,-0.5);
\draw (6.5,6.5) -- (6.5,9.5);


\foreach \x in {0,...,9}
{
  \foreach \y in {0,...,9}
  {
    \node[place1] () at (\x,\y) {};
  }
}

\foreach \x in {3,...,6}
{
  \node[place2] () at (\x,2) {};
  \node[place2] () at (\x,7) {};
}
\foreach \y in {3,...,6}
{
  \node[place2] () at (2,\y) {};
  \node[place2] () at (7,\y) {};
}


\foreach \x in {3,...,6}
{
  \node[place3] () at (\x,1) {};
  \node[place3] () at (\x,8) {};
}
\foreach \y in {3,...,6}
{
  \node[place3] () at (1,\y) {};
  \node[place3] () at (8,\y) {};
}


\node[place3] () at (2,2) {};
\node[place3] () at (7,2) {};
\node[place3] () at (2,7) {};
\node[place3] () at (7,7) {};


\end{tikzpicture}
}
\end{figure}

Pour appliquer cette technique, il faut toutefois savoir à l'avance combien de fois le stencil sera itéré, cela est possible avec la syntaxe proposée pour le chaînage de stencils, mais cela n'est pas possible pour le cas général où une nouvelle syntaxe est alors nécessaire (par exemple en intercalant le nombre d'itérations dans la syntaxe présentée dans le code \ref{lst:agreg1_stencil}). 

Une nouvelles fois c'est la spécialisation partielle qui permettrait de mettre en œuvre le déclenchement automatique de l'utilisation de cette technique. C'est aussi par la spécialisation partielle que pourrait être choisie la forme la plus adaptée des blocs pour le tuilage. 

%% Chap4:
%% -> mise en \oe uvre, éléments d'implantation
%%    i. e. C++, SYCL

%% -*- eval: (flyspell-mode 1); -*-

\chapter{\'Evaluation de la bibliothèque}


\section{\'Evaluation qualitative}

Comparaisons d'expressivités (nombre de lignes, nombres de symboles, typage à la compilation), factorisation de code
Reprendre slides : différents cas de tests


\section{Performances quantitatives sur machines dédiées}

Description du matériel

Mesures de temps

Mesures de compteurs

Vérification des valeurs numériques

%% Chap5:
%% -> évaluation, méthodologie présentation résultat

%% -*- eval: (flyspell-mode 1); -*-

\chapter{Conclusion}
%\addcontentsline{toc}{chapter}{Conclusion}

Au cours du stage a été développée un DSEL permettant de décrire facilement les calculs à base de stencils et d'obtenir des temps d'exécution en corrélation avec les capacités actuelles des machines. Il s'agit là de deux problématiques primordiales du domaine du calcul haut performance : comment faciliter la programmation de codes efficaces par un utilisateur quelconque, et comment les rendre rapides à l'exécution. Ce DSEL utilise les moyens les plus modernes disponibles tels que le \textsf{C++14} ou \textsf{SYCL} et répond au manque d'outils adaptés pour les stencils. L'évaluation du DSEL a permis d'établir que des améliorations restent à prévoir, elles seront toutefois facilités par les techniques qu'il utilise déjà. 

\section{Bilan}

Les quelques langages (tels que \textsf{Blitz++} ou \textsf{Pochoir}) existants pour les calculs de stencils ont de nombreux désavantages à la fois dans leur syntaxe et aussi dans leur performances car peu sont adaptés aux machines parallèles hétérogènes actuelles. Après une analyse des éléments nécessaires à la description d'un stencil, des objectifs ont été énoncés concernant un futur outil dédié aux stencils. Ils concernent tout autant l'expressivité de l'outil --- le respect de la sémantique des stencils ---, que les performances qu'il doit apporter --- notamment par la technique de tuilage. Ces objectifs ont pour la plupart été repris lors de l'implantation du DSEL basé sur \textsf{triSYCL} qui est opérationnel et a été testé sur un exemple de problème courant : Jacobi en deux dimensions. 

D'autres choix auraient été possibles concernant la conception ainsi que l'implantation de l'outil. Concernant la conception il n'était pas obligatoire par exemple de se focaliser sur un DSEL : une extension de compilateur permettant de détecter automatiquement les stencils dans le code serait certainement très bénéfique pour l'utilisateur bien que difficile à réaliser. Quant à l'implantation, le DSEL aurait pu utiliser un autre langage tel que \textsf{Python}, très apprécié grâce à sa syntaxe pour les tableaux. D'autres bibliothèques de parallélisation adaptées aux machines hétérogènes auraient aussi pu être utilisées par le DSEL au lieu de \textsf{triSYCL} : \textsf{SkePU} qui s'appuie sur \textsf{StarPU}, ou directement \textsf{StarPU}.

Les tests effectués ont toutefois démontré que notre approche était viable tant pour l'expressivité que pour les performances offertes par le DSEL. Il s'agit donc maintenant de le compléter afin de l'intégrer dans des applications existantes telles que \textsf{QIRAL}.

\section{Perspectives}

À court terme les évolutions pour le DSEL concernent le nombre de dimensions autorisées pour les stencils (uniquement deux actuellement), quelques éléments de syntaxe (pour la gestion de types complexes d'éléments telles que des matrices) et le choix automatique de l'implantation optimale (technique de tuilage ou non).

Par ailleurs la représentation en mémoire des données utilisées par les stencils mérite une attention poussée qui pourrait conduire à de nouvelles implantations, de même que la gestion des aspects distribués (si plusieurs machines distinctes connectées entre elles). Les caractéristiques des  processeurs (CPU, GPU ou \emph{many integrated core}) doivent aussi permettre d'adapter le code produit par le DSEL, et il est par exemple nécessaire d'améliorer et uniformiser la façon dont ces caractéristiques sont récupérées.

Enfin à long terme le DSEL pourrait s'étendre à d'autres types de calculs que les stencils, éventuellement en reprenant le formalisme des squelettes de calculs dont les \emph{map/reduce} qui sont très courants dans les problèmes actuels.

Les perspectives à court terme seront atteintes prochainement et le DSEL produit sera proposé avec les codes d'exemples de la bibliothèque OpenSource \textsf{triSYCL} développée par Ronan Keryell. Cela permettra notamment à la communauté elle-même de jauger et améliorer directement ce nouvel outil. 


%% Conclusion:
%% -> ce qui a été fait (résumé)
%% -> ce qui aurait pu être fait (qqc de crédible, choix alternatifs par exemple SYCLONE, SkePU, StarPU)
%% -> ce qui reste à faire (à court (plutôt technique), moyen et long terme) (aspect ditribué ? pour aute algo que stencils comme map/reduce ?)


\phantomsection
\chapter*{Bibliographie}
\addcontentsline{toc}{chapter}{Bibliographie}

%%\nocite{Web3}

\section*{Articles}

\printbibliography[heading=none,nottype=thesis,nottype=online]

\section*{Thèses}

\printbibliography[heading=none,type=thesis]

\section*{Autres}

\printbibliography[heading=none,type=online]

\end{document}
